% ============================================================
% Capitolo 8 — Model Levels
% ============================================================

\chapter{Model Levels: quale AI per quale task}
\label{ch:model-levels}

C'è un errore comune nell'uso degli agenti AI: usare sempre il modello più potente disponibile. Il primo motivo è il costo: i modelli più potenti costano di più in token. Il secondo è la latency: su task semplici e frequenti la differenza è percepibile e non giustificata.

ADLC introduce i Model Levels per risolvere questo: una scala da 1 a 7 che mappa la complessità del task al modello appropriato.

\section{La scala 1--7}
\label{sec:08-scala}

\begin{table}[htbp]
  \centering
  \begin{tabular}{cll}
    \toprule
    \textbf{Level} & \textbf{Token range (totali)} & \textbf{A cosa serve} \\
    \midrule
    1 & $<$ 4k & Piccoli edit, docs, formattazione \\
    2 & 4k--8k & Modifiche localizzate, test semplici \\
    3 & 8k--16k & Implementazione standard, debugging \\
    4 & 16k--32k & Feature multi-file, design moderato \\
    5 & 32k--64k & Refactor complessi, sicurezza \\
    6 & 64k--120k & Architettura, debugging profondo \\
    7 & $>$ 120k & Mission-critical, alta ambiguità \\
    \bottomrule
  \end{tabular}
  \caption{Scala dei Model Level ADLC}
  \label{tab:08-scala}
\end{table}

Il mapping ai modelli concreti (Anthropic, OpenAI, Gemini) vive in \texttt{.adlc/manifest.json\#model\_levels} e si legge con:

\begin{lstlisting}[language=,caption={Visualizzazione mapping modelli}]
bash .adlc/tools/show-models.sh   # Linux/macOS/WSL
.\.adlc\tools\show-models.ps1     # Windows PowerShell
\end{lstlisting}

\section{Come stimare il livello di un task}
\label{sec:08-stima}

\subsection*{Passo 1 --- Stima dei token}

La stima è \texttt{input / output / totale}. Alcune euristiche pratiche:

\begin{table}[htbp]
  \centering
  \begin{tabular}{ll}
    \toprule
    \textbf{Input} & \textbf{Token approssimativi} \\
    \midrule
    File da 100 righe & 500--800 \\
    \texttt{\_CONTEXT.md} compilato & 500--1.000 \\
    3 file da 200 righe & 3.000--5.000 \\
    Modulo da 10 file & 10.000--20.000 \\
    \bottomrule
  \end{tabular}
  \caption{Euristiche per la stima dei token in input}
  \label{tab:08-euristiche}
\end{table}

Esempio per TaskFlow API --- Lorenzo implementa \texttt{GET /tasks} con filtri:
\begin{itemize}
  \item Input: \texttt{\_CONTEXT.md} (700) + route esistente (600) + service (400) + schema (300) = $\sim$2.000
  \item Output: nuovo endpoint + test $\sim$ 2.500
  \item Totale: $\sim$4.500 token $\to$ range 4k--8k $\to$ \textbf{Livello 2} (dalla tabella~\ref{tab:08-scala})
\end{itemize}

\subsection*{Passo 2 --- Applica il risk floor}

Il task è classificato MEDIUM. Il risk floor per MEDIUM è livello 3. Anche se la stima token suggerisce livello 2, il level sale a 3.

\begin{lstlisting}[language=,caption={AI Sizing in \_CONTEXT.md}]
| Active Task Token Est.   | 2000/2500/4500  |
| Active Task Model Level  | 3 Sonnet Low    |
\end{lstlisting}

\section{I risk floor in dettaglio}
\label{sec:08-risk-floor}

\begin{table}[htbp]
  \centering
  \begin{tabular}{ll}
    \toprule
    \textbf{Condizione} & \textbf{Minimum Level} \\
    \midrule
    Task MEDIUM & 3 \\
    Task HIGH & 5 \\
    Auth, secrets, compliance, produzione & 6 \\
    Architettura o cambiamenti cross-modulo & 6 \\
    CRITICAL & 7 \\
    \bottomrule
  \end{tabular}
  \caption{Risk floor per condizione}
  \label{tab:08-risk-floor}
\end{table}

Il risk floor è cumulativo: un task HIGH (floor 5) che tocca auth (floor 6) usa il massimo --- livello 6.

\section*{\LabelSummary}

\begin{itemize}
  \item I Model Level (1--7) mappano la complessità del task al modello AI appropriato.
  \item La stima è \texttt{input/output/totale} in token; il mapping vendor concreto è in \texttt{manifest.json}.
  \item Il risk floor forza un minimum level indipendentemente dalla stima token.
  \item La stima migliora con la pratica; le prime settimane è normale sbagliare.
\end{itemize}
